\chapter{المقاطعات وبرامج التشغيل}
\label{chap:interrupt}

يتولى \indextext{برنامج التشغيل (Driver)} إدارة أحد ملحقات العتاد الفيزيائي بداخل النواة (مثل وحدة UART لتوصيل لوحة المفاتيح والأنبوب النصفي، أو محرك الأقراص).
ينظم برنامج التشغيل إرسال الأوامر إلى الجهاز واستقبال البيانات منه عبر المقاطعات وإشارات العتاد.

تصل \indextext{مقاطعات الأجهزة (Device Interrupts)} بصورة غير متزامنة مع تنفيذ التعليمات؛
حيث يرسل موصل المقاطعات بـ RISC-V والمسمى \indextext{PLIC (Platform-Level Interrupt Controller)} إشارة المقاطعة للمعالج عند وصول بايت جديد من اللوحة أو اكتمال القراءة من القرص.

تقوم النواة بمعالجة المقاطعة عبر الدالة \texttt{devintr()}، التي تتعرف على مصدر المقاطعة وتستدعي برنامج التشغيل المناسب (مثل \texttt{uartintr()} أو \texttt{virtio_disk_intr()}).
يقوم برنامج التشغيل بنقل البيانات من سجلات الجهاز المربوطة بالذاكرة (MMIO) إلى ذاكرة النواة، ثم يوقظ العمليات المنتظرة لتلك البيانات.
